% Exact main-gate reduction, work note, 19 November2026.
% The final polynomial-kernel inclusion is explicitly uncomputed.
\subsection*{From a full profinite nonperipheral normalizer to the single-line matrix gate}

Put $X_0=E_0\setminus\{O\}$, with
$E_0:y^2=x^3-x$ over $\Q$, and use a rational tangential
basepoint at $O$. Write $F=\pi_1(X_{0,\overline\Q})$ for
the full profinite geometric group and
$\Pi=F\rtimes s_O(G_\Q)$ for its arithmetic group.
Let $A=\overline{\langle z\rangle}\subset F$ be nontrivial
and procyclic, and let its arithmetic normalizer $N$
have open image in $G_{\Q(i)}$. This image is also open
in $G_\Q$. We construct the exact image of every
nonperipheral such $z$ in the single-line problem.

\paragraph{1. The original full-class proper-homology detector.}
Apply Nakamura1994, Lemma~2.1.4, printed pp.~96--97,
to the full class $\mathcal C$ of all finite groups\cite{AuditNakamura1994}.
Its hypotheses are $n\geq1$ and $2-2g-n<0$, and they
hold for $(g,n)=(1,1)$. Let $Z$ be the union of the
conjugates of all cusp inertia subgroups, exactly as
defined in that lemma. For $z\notin Z$, the lemma gives
a prime $\ell$ and an open subgroup $H\subset F$
containing $z$ such that
\[
             z\notin [H,H]H^\ell\langle Z\cap H\rangle .    \tag{FC.1}
\]
All products and generated subgroups here have the
closed meanings used in the profinite statement.
Because $H$ is closed and contains $z$, it contains
the entire $A=\overline{\langle z\rangle}$.

Let $Y_H$ be the actual geometric finite etale cover
associated with $H$, and $C_H$ its smooth proper
compactification. The elements of $Z\cap H$ are
precisely the inertia elements over the boundary
of this cover. Killing their closed normal closure
is the actual filling map
$H\twoheadrightarrow\pi_1(C_H)$. Abelianizing modulo
$\ell$ identifies the quotient in (FC.1) with
$H_1(C_H,\mathbf F_\ell)$. Consequently the actual
class $c_z$ of $z$ in
\[
                    H_1(C_H,\Z_\ell)=T_\ell J_{C_H}
\]
reduces to a nonzero class modulo $\ell$.
This is the image of the specified element $z$,
not a freely chosen lift of an unrelated residue
class. Thus $c_z\ne0$. The module is free over
$\Z_\ell$, so its image in
$H_1(C_H,\Q_\ell)$ is also nonzero.
In particular $z$ has infinite order.

The procyclic group $A$ has a canonical product of
its prime-primary factors. Its homomorphism to
$H_1(C_H,\Z_\ell)$ kills every factor at primes
different from $\ell$. Let $u$ be the $\ell$-component
of $z$. Then its actual class equals $c_z$, so
$A_\ell\simeq\Z_\ell$ and $[u]\ne0$ in proper
rational homology. The subgroup $A_\ell$ is
characteristic in $A$, hence normalized by $N$.
No assertion about the other prime components
or a common cyclotomic character has been used.

\paragraph{2. A rational normal core of genus at least three.}
Fix the actual finite etale cover of $X_0$ supplied
by F.3.16 whose proper compactification is the
Fermat quartic of genus three. Let $H_3\subset F$
be its geometric subgroup. This geometric cover
need only be defined over a finite extension;
its genus and its actual morphism are sufficient.
Set $H_0=H\cap H_3$ and form
\[
                  S=\bigcap_{\gamma\in\Pi}
                          \gamma H_0\gamma^{-1}.           \tag{FC.2}
\]
There are only finitely many open subgroups of a
fixed index in the finitely generated profinite
group $F$: a continuous permutation representation
is determined by the images of its fixed finite
generating set. Therefore (FC.2) is an open
subgroup, normal in $\Pi$, contained in both
$H$ and $H_3$.

For some integer $r\geq0$, the element
$u'=u^{\ell^r}$ generates $A_\ell\cap S$.
The actual inclusion $S\hookrightarrow H$ and
the finite map of compactifications give
\[
 H_1(C_S,\Q_\ell)\longrightarrow H_1(C_H,\Q_\ell),
                    \qquad [u']\longmapsto\ell^r[u].       \tag{FC.3}
\]
The right side is nonzero. Thus the proper class
of $u'$ on the source is nonzero. The same cover
has a finite nonconstant map of proper curves
$C_S\to C_3$ because $S\subset H_3$.
Riemann--Hurwitz in characteristic zero gives
$2g(C_S)-2\geq\deg(C_S/C_3)(2\cdot3-2)$;
in particular $g(C_S)\geq3$.

The arithmetic subgroup
\[
                     \Pi_Y=S\rtimes s_O(G_\Q)
\]
defines an actual geometrically connected cover
$Y/\Q$ of $X_0$. Its chosen cusp decomposition
subgroup is
\[
 \Pi_Y\cap(I_O\rtimes s_O(G_\Q))
                  =(S\cap I_O)\rtimes s_O(G_\Q).
\]
Its arithmetic residue quotient is $G_\Q$, so
the corresponding point $y$ of the smooth proper
compactification $C=\overline Y$ is rational.
This is exactly the tangential descent map
proved in F.3.23. We have $g(C)\geq3$.

The subgroup $N'=N\cap\Pi_Y$ has finite index in
$N$, hence open arithmetic image. It normalizes
$A_\ell\cap S$: both $A_\ell$ and $S$ are preserved
by $N$. Its geometric kernel lies in $S$, acts
inner on $S$, and therefore acts trivially on
the proper homology of $C$. The nonzero line
$\Q_\ell[u']$ is consequently an actual
arithmetic normalizer line in $T_\ell J_C\otimes\Q_\ell$.

\paragraph{3. Filling to an actual once-punctured target.}
The open immersion
\[
                         Y\hookrightarrow X=C\setminus\{y\}
                                                               \tag{FC.4}
\]
is defined over $\Q$. Its fundamental-group map
is a geometric surjection and induces the
identity on the arithmetic quotient. The
proper homology map is the identity on
$H_1(C,\Q_\ell)$, so the class in (FC.3)
remains nonzero. Compose the geometric map
with the maximal pro-$\ell$ quotient
\[
                     S\longrightarrow
                     P=\pi_1(X_{\overline\Q})^{(\ell)}.
                                                               \tag{FC.5}
\]
The image of $A_\ell\cap S$ is a nontrivial
closed procyclic subgroup with nonzero
abelianized class. In fact it is isomorphic
to $\Z_\ell$: a nontrivial closed kernel of
a homomorphism from $\Z_\ell$ would have
finite quotient, whose rational homology
would be zero. The image of $N'$ in the
arithmetic geometrically pro-$\ell$ group
normalizes this subgroup and has the same
open arithmetic image. Compactness makes
that image closed.

The curve $X$ has a rational tangential
basepoint, genus at least three, and one
puncture. Its degree-one homology is the
proper homology of $C$. Thus (FC.1)--(FC.5)
give the exact composite from a full
profinite nonperipheral normalizer to the
single nonzero homology-line configuration
on one actual once-punctured rational curve.
The common kernel of all pro-prime quotients
of the original $F$ is not quotiented away
before applying the detector: a prime is
chosen only after passing to $H$.

\paragraph{The actual nonzero arithmetic derivation.}
Choose a finite number field $L$ whose absolute Galois
group lies in the open arithmetic image of the
normalizer in (FC.5). Nakamura1996, printed
pp.~115--117\cite{AuditNakamura1996}, gives the actual nonzero degree-six
quotient for this original curve $X$: the rank
sandwich in Corollary~4.2 has both endpoints equal
to one in degree six, Lemma~4.5 gives the actual
fibre surjection and finite-base invariance, and
(4.7.1)--(4.7.2) identify the geometric kernel and
the equivariant universal quotient. Thus there
exists an actual
\[
 \sigma\in G_L(6),\qquad D_\sigma\in U_{g,6},
                           \qquad q_{g,6}(D_\sigma)\ne0.
                                                               \tag{FC.5a}
\]
Lift this same $\sigma$ to the image of $N'$.
Its action on the nonzero proper homology class
$a$ of the procyclic generator is the identity,
so its exponent on that procyclic subgroup is
one. In the rational completed free Lie algebra
the normalizing lift and the canonical special
representative differ by an inner conjugation:
\[
       \phi_\sigma(\log u')
                    =\exp(\operatorname{ad}Z)(\log u').
\]
The degree-one part of $Z$ is a multiple of $a$.
Remove it by a rational power of the full
$\log u'$, which commutes with itself. The
degree-two through degree-five parts then
centralize $a$ and vanish: the associative
centralizer is $\Q_\ell[a]$, whose primitive
part is $\Q_\ell a$. Degree seven gives the
actual equality
\[
                         D_\sigma(a)=[Z_6,a].               \tag{FC.5b}
\]
This is the full-logarithm correction proved
in F.3.29, on the present genus-$g$ free group.
Apply the actual Levi invariant projector
of F.3.30. The vector
$D=e_S D_\sigma$ remains in $U_{g,6}\cap P_a$,
is Levi invariant, and has
$q_{g,6}(D)=q_{g,6}(D_\sigma)\ne0$.
For $g=3$ this contradicts F.3.29 already;
every remaining configuration therefore has
$g\geq4$ and enters the invariant matrix
below with a nonzero actual quotient.

\paragraph{4. The finite scalar-polynomial gate.}
F.3.29 handles genus three. In each genus
$g\geq4$, put $s=g-1$. The exact construction
of F.3.30 defines one finite integer
polynomial matrix $M(s)$ on its five
pure-complement inputs and215 raw Dynkin
inputs. Its rows include the constant
suffix-Lie and $b$-leading consistency
equations and all scalar invariant-dual
contractions of the coupled mixed map
\[
 D\longmapsto\bigl(\delta\mapsto
    (\operatorname{ES}_7([D,\delta]),R([D,\delta]))\bigr).
\]
Every row is obtained by the specified
finite pairing diagrams; every closed
$W$-loop contributes a factor $2s$ with
its orientation sign. There are finitely
many such dual diagrams in the fixed
tensor degrees. More explicitly, curry the
map on $\Lambda^3V$ by alternating the three
dual input slots. For the ES output use
the cyclic average in its seven output
slots before applying a dual pairing
diagram. These are literal maps on the
exterior input and the cyclic-coinvariant
output, not functionals on a freely
chosen representative. In characteristic
zero the invariant dual tensors separate
these actual quotient targets. Denominators in (LI.2)
are fixed integers and can be cleared.
The output map $K_{\mathcal D}(s)$ is
likewise an explicitly defined polynomial
matrix of five complete contractions.

Let $L$ be a fixed three-row rational matrix
whose kernel is the span of the two
constant detector vectors of $Z_1,Z_2$.
Their determinant-$10752$ minor gives, in the raw
contraction convention, the explicit choice
\[
 L=\begin{pmatrix}-1&4&0&0&0\\1&0&1&0&0\\0&0&0&1&1\end{pmatrix}.
\]
Its kernel is exactly the span of the two displayed
columns in F.3.29. This also fixes the matching signs
in the scalar output quotient.
The remaining, unverified finite algebraic
assertion is
\[
 \ker M(s)\subset\ker\bigl(LK_{\mathcal D}(s)\bigr)
                    \qquad(s=3,4,5,\ldots).                \tag{FC.6}
\]
This is a sufficient scalar-image
certificate for excluding the actual
configurations in (FC.5), by the actual
projector and five-detector map of F.3.30.
It is not stated to be a necessary
criterion for the full subgroup assertion.

Precisely, every nonperipheral full
profinite subgroup under consideration
produces, by (FC.1)--(FC.5b), an integer
$s\geq3$ and a parameter vector $v$ over
$\Q_\ell$ with
\[
               M(s)v=0,\qquad LK_{\mathcal D}(s)v\ne0.
                                                               \tag{FC.7}
\]
The first equality is the actual coupled
geometric-normalizer equation. The
inequality follows from the proved
five-detector inclusion (LI.6) and the
nonzero quotient in (FC.5a). Thus this
is an unconditional witness map from
full subgroup counterexamples to failures
of (FC.6). Its converse is not asserted.
Since the matrices at an integer $s$
have rational entries, existence of the
linear-algebra witness over $\Q_\ell$
is equivalent to failure of the same
kernel inclusion over $\Q$.

The assertion (FC.6) is finitely decidable
once the finite matrices have been formed.
Compute Smith normal form over $\Q[s]$.
Away from the zeros of its nonzero invariant
factors, the generic kernel has an explicit
polynomial basis after the unimodular
column transformation; check that
$LK_{\mathcal D}$ kills that basis as
polynomial identities. The finitely many
integer roots $s\geq3$ of the nonzero
invariant factors are then separate exact
rational matrix tests. This procedure also
covers rank-zero matrices. The finite
integer roots can be determined after
clearing denominators by factorization or
the rational-root theorem. No sampling
or polynomial interpolation alone is
used to justify the quantified assertion.

The mixed matrix has not yet been
computed, and (FC.6) has not been verified.
The proven reduction exhibits how that
single exact finite gate bears on the
full profinite subgroup, including the
residual pronilpotent kernel. It does not
declare the gate passed or the original
LM6 A.22 assertion resolved.
